\documentclass[10pt,a4paper]{amsart}
\usepackage{amsmath,amssymb,amsthm,graphicx}
% Retained minimal preamble for the Stacks Project chapter varieties.tex.
% Verbatim excerpt of the chapter preamble: Xy-pic setup (native preamble lines 14-18), the native theorem
% declarations with the mathematical macros (native preamble lines 154-217), and the project own no-op
% definitions of the reference/slogan/history display environments (native wrapper.tex lines 1-3) that the
% counted statement uses for its slogan. Copied mechanically; nothing added, removed or redefined.
% For commutative diagrams we use Xy-pic
\usepackage[all]{xy}

% We use 2cell for 2-commutative diagrams.
\xyoption{2cell}
% Theorem environments.
%
\theoremstyle{plain}
\newtheorem{theorem}[subsection]{Theorem}
\newtheorem{proposition}[subsection]{Proposition}
\newtheorem{lemma}[subsection]{Lemma}

\theoremstyle{definition}
\newtheorem{definition}[subsection]{Definition}
\newtheorem{example}[subsection]{Example}
\newtheorem{exercise}[subsection]{Exercise}
\newtheorem{situation}[subsection]{Situation}

\theoremstyle{remark}
\newtheorem{remark}[subsection]{Remark}
\newtheorem{remarks}[subsection]{Remarks}

\numberwithin{equation}{subsection}

% Macros
%
\def\lim{\mathop{\mathrm{lim}}\nolimits}
\def\colim{\mathop{\mathrm{colim}}\nolimits}
\def\Spec{\mathop{\mathrm{Spec}}}
\def\Hom{\mathop{\mathrm{Hom}}\nolimits}
\def\Ext{\mathop{\mathrm{Ext}}\nolimits}
\def\SheafHom{\mathop{\mathcal{H}\!\mathit{om}}\nolimits}
\def\SheafExt{\mathop{\mathcal{E}\!\mathit{xt}}\nolimits}
\def\Sch{\mathit{Sch}}
\def\Mor{\mathop{\mathrm{Mor}}\nolimits}
\def\Ob{\mathop{\mathrm{Ob}}\nolimits}
\def\Sh{\mathop{\mathit{Sh}}\nolimits}
\def\NL{\mathop{N\!L}\nolimits}
\def\CH{\mathop{\mathrm{CH}}\nolimits}
\def\proetale{{pro\text{-}\acute{e}tale}}
\def\etale{{\acute{e}tale}}
\def\QCoh{\mathit{QCoh}}
\def\Ker{\mathop{\mathrm{Ker}}}
\def\Im{\mathop{\mathrm{Im}}}
\def\Coker{\mathop{\mathrm{Coker}}}
\def\Coim{\mathop{\mathrm{Coim}}}

% Boxtimes
%
\DeclareMathSymbol{\boxtimes}{\mathbin}{AMSa}{"02}

%
% Macros for moduli stacks/spaces
%
\def\QCohstack{\mathcal{QC}\!\mathit{oh}}
\def\Cohstack{\mathcal{C}\!\mathit{oh}}
\def\Spacesstack{\mathcal{S}\!\mathit{paces}}
\def\Quotfunctor{\mathrm{Quot}}
\def\Hilbfunctor{\mathrm{Hilb}}
\def\Curvesstack{\mathcal{C}\!\mathit{urves}}
\def\Polarizedstack{\mathcal{P}\!\mathit{olarized}}
\def\Complexesstack{\mathcal{C}\!\mathit{omplexes}}
% \Pic is the operator that assigns to X its picard group, usage \Pic(X)
% \Picardstack_{X/B} denotes the Picard stack of X over B
% \Picardfunctor_{X/B} denotes the Picard functor of X over B
\def\Pic{\mathop{\mathrm{Pic}}\nolimits}
\def\Picardstack{\mathcal{P}\!\mathit{ic}}
\def\Picardfunctor{\mathrm{Pic}}
\def\Deformationcategory{\mathcal{D}\!\mathit{ef}}
\newenvironment{reference}{}{}
\newenvironment{slogan}{}{}
\newenvironment{history}{}{}
\usepackage{cleveref}
\crefname{originallemma}{original lemma}{original lemmas}
\Crefname{originallemma}{Original lemma}{Original lemmas}
\makeatletter
\newlabel{algebra-lemma-local-flat-ff}{{00HR}{source}{}{sourcepointer-1}{}}
\newlabel{algebra-lemma-local-flat-ff@cref}{{[originallemma][0][]00HR}{[1][0][]source}{}{}{}}
\makeatother
\crefname{originallemma}{original lemma}{original lemmas}
\Crefname{originallemma}{Original lemma}{Original lemmas}
\makeatletter
\newlabel{schemes-lemma-points-fibre-product}{{01JT}{source}{}{sourcepointer-2}{}}
\newlabel{schemes-lemma-points-fibre-product@cref}{{[originallemma][0][]01JT}{[1][0][]source}{}{}{}}
\makeatother
\title{Local rings under a ground field extension: faithful flatness and the residue field at a point of the base change}
\author{The Stacks Project Authors}
\date{stacks-project:varieties.tex:local-snapshot-20261011; GNU FDL 1.2 or later}
\begin{document}
\maketitle
\noindent\textit{Editorial note.} All delivered mathematical text is the author's own, copied from the pinned Stacks Project chapter file varieties.tex; no statement or proof step was rewritten, reordered or removed.\par\medskip
\noindent\textit{Editorial note.} The proof contains the parenthetical '(details omitted)' at the identification of the prime used for the localization; the parenthetical is reproduced verbatim. The surrounding factorization and the description of the prime are written out, so the omitted check is bounded.\par\medskip
\noindent\textit{Editorial note.} The counted spans carry 2 navigation stubs and no bibliography citation.\par\medskip
\noindent\textit{Editorial note.} The preamble is the chapter's retained minimal preamble (sources/stacks-varieties/retained-preamble.tex lines 1-77), a verbatim excerpt of the native preamble.\par\medskip
\noindent\textit{Editorial note.} The label lemma-change-fields-flat is defined only in varieties.tex, so no cross-chapter restatement had to be weighed for this card.\par\medskip
\noindent\textit{Allowed standard prerequisites (editorial).}\begin{itemize}
\item A local flat homomorphism of local rings with a residue field extension and a nilpotent maximal ideal structure is faithfully flat. Algebra, Lemma algebra-lemma-local-flat-ff, invoked at proof lines 258-259. Supplies faithful flatness in part (1).
\item Points of a fibre product and their residue fields. Schemes, Lemma schemes-lemma-points-fibre-product, invoked at proof lines 260-261. Supplies the description of the residue field at a point of the base change used in part (2).
\item Exactness of localization and of tensoring with a field extension. Standard, used at proof lines 279-280. Supplies part (4) from part (3).
\end{itemize}\medskip
\ifdefined\hypertarget\hypertarget{sourcepointer-1}{}\fi
\noindent\textit{Source reference.} Original lemma 00HR: sources/stacks-chow/tags.txt, source lines 653--653. Pointer outside the counted span, invoked once for context; the load-bearing argument of the counted statement is the author's proof reproduced in full in the proof span.\par\medskip
\ifdefined\hypertarget\hypertarget{sourcepointer-2}{}\fi
\noindent\textit{Source reference.} Original lemma 01JT: sources/stacks-chow/tags.txt, source lines 2026--2026. Pointer outside the counted span, invoked once for context; the load-bearing argument of the counted statement is the author's proof reproduced in full in the proof span.\par\medskip
\begin{lemma}
\label{lemma-change-fields-flat}
Let $K/k$ be an extension of fields. Let $X$ be scheme over $k$
and set $Y = X_K$. If $y \in Y$ with image $x \in X$, then
\begin{enumerate}
\item $\mathcal{O}_{X, x} \to \mathcal{O}_{Y, y}$ is a
faithfully flat local ring homomorphism,
\item with $\mathfrak p_0 = \Ker(\kappa(x) \otimes_k K \to \kappa(y))$
we have $\kappa(y) = \kappa(\mathfrak p_0)$,
\item $\mathcal{O}_{Y, y} = (\mathcal{O}_{X, x} \otimes_k K)_\mathfrak p$
where $\mathfrak p \subset \mathcal{O}_{X, x} \otimes_k K$ is the inverse
image of $\mathfrak p_0$.
\item we have
$\mathcal{O}_{Y, y}/\mathfrak m_x\mathcal{O}_{Y, y} =
(\kappa(x) \otimes_k K)_{\mathfrak p_0}$
\end{enumerate}
\end{lemma}

\begin{proof}
We may assume $X = \Spec(A)$ is affine. Then $Y = \Spec(A \otimes_k K)$.
Since $K$ is flat over $k$, we see that $A \to A \otimes_k K$ is flat.
Hence $Y \to X$ is flat and we get the first statement if we also
use Algebra, Lemma \ref{algebra-lemma-local-flat-ff}.
The second statement follows from
Schemes, Lemma \ref{schemes-lemma-points-fibre-product}.
Now $y$ corresponds to a prime ideal $\mathfrak q \subset A \otimes_k K$
and $x$ to $\mathfrak r = A \cap \mathfrak q$. Then $\mathfrak p_0$
is the kernel of the induced map
$\kappa(\mathfrak r) \otimes_k K \to \kappa(\mathfrak q)$.
The map on local rings is
$$
A_\mathfrak r \longrightarrow (A \otimes_k K)_\mathfrak q
$$
We can factor this map through
$A_\mathfrak r \otimes_k K = (A \otimes_k K)_{\mathfrak r}$
to get
$$
A_\mathfrak r \longrightarrow A_\mathfrak r \otimes_k K
\longrightarrow (A \otimes_k K)_\mathfrak q
$$
and then the second arrow is a localization at some prime. This prime ideal
is the inverse image of $\mathfrak p_0$ (details omitted) and this
proves (3). To see (4) use (3) and that localization and $- \otimes_k K$
are exact functors.
\end{proof}

\end{document}
